\begin{examplebox}
	\textbf{\textcolor{CExample}{例 1（基础）}}

	\textbf{\textcolor{CExample}{题目：}}
	求函数 $f(x)=3x+\frac{12}{x}$（$x>0$）的最小值，并求出相应的 $x$ 值；再求 $x<0$ 时 $f(x)$ 的最大值及相应 $x$ 值。

	\textbf{\textcolor{CExample}{解题步骤：}}
	\begin{description}[style=nextline, leftmargin=2.8em, labelsep=0.5em, itemsep=0.45em]
		\item[\textbf{第一步（识别结构）：}] 确定函数为对勾函数标准形式 $3x+\frac{12}{x}$，其中 $A=3$, $B=12$, 且 $A,B>0$.
			\[
				f(x)=3x+\frac{12}{x}
			\]
			\par\textbf{理由：} 标准对勾结构可直接套用结论.

		\item[\textbf{第二步（正数最小值）：}] 当 $x>0$ 时，由结论
			\[
				\begin{aligned}
					f(x) &\ge 2\sqrt{AB} \\
					&= 2\sqrt{3\cdot12} \\
					&= 12,
			\end{aligned}
			\]
			等号成立当且仅当 $x=2$.
			\par\textbf{理由：} 直接应用 $x>0$ 时 $f(x)\ge 2\sqrt{AB}$ 公式.

		\item[\textbf{第三步（负数最大值）：}] 当 $x<0$ 时，由结论
			\[
				\begin{aligned}
					f(x) &\le -2\sqrt{AB} \\
					&= -12,
			\end{aligned}
			\]
			等号成立当且仅当 $x=-2$.
			\par\textbf{理由：} 直接应用 $x<0$ 时 $f(x)\le -2\sqrt{AB}$ 公式.
	\end{description}

	\textbf{\textcolor{CExample}{关键结论：}}
	\textbf{最小值为 $12$（$x=2$）；最大值为 $-12$（$x=-2$）}.

	\medskip
	\textbf{\textcolor{CExample}{例 2（稍变形）}}

	\textbf{\textcolor{CExample}{题目：}}
	求函数 $y=\frac{2x^2+3x+3}{x+1}$（$x\neq-1$）的值域.

	\textbf{\textcolor{CExample}{解题步骤：}}
	\begin{description}[style=nextline, leftmargin=2.8em, labelsep=0.5em, itemsep=0.45em]
		\item[\textbf{第一步（除法变形）：}] 将分子除以分母，得
			\[
				y = 2x+1+\frac{2}{x+1}.
			\]
			\par\textbf{理由：} 将分式化为一次式加余式的形式，便于换元.

		\item[\textbf{第二步（换元转化）：}] 令 $t=x+1$，则 $x=t-1$，代入得
			\[
				\begin{aligned}
					y &= 2(t-1)+1+\frac{2}{t} \\
					&= 2t+\frac{2}{t} -1,
			\end{aligned}
			\]
			其中 $t\in(-\infty,0)\cup(0,+\infty)$.
			\par\textbf{理由：} 通过线性换元将函数化为 $2t+\frac{2}{t}-1$，即 $At+\frac{B}{t}+C$ 形式.

		\item[\textbf{第三步（结论求值域）：}] 对 $g(t)=2t+\frac{2}{t}$（标准对勾，$A=2,B=2>0$），由结论知 $g(t)$ 的值域为 $(-\infty,-4]\cup[4,+\infty)$。于是 $y=g(t)-1$ 的值域为
			\[
				(-\infty,-5]\cup[3,+\infty).
			\]
			\par\textbf{理由：} 先求对勾部分的值域，再整体减 $1$（常数平移）即得原函数值域.
	\end{description}

	\textbf{\textcolor{CExample}{关键结论：}}
	\textbf{值域为 $(-\infty,-5]\cup[3,+\infty)$}.

	\medskip
	\textbf{\textcolor{CExample}{例 3（综合应用）}}

	\textbf{\textcolor{CExample}{题目：}}
	求函数 $y=\frac{2x^2+3x+3}{x+1}$ 在 $x>-1$ 时的最小值，并求相应的 $x$ 值.

	\textbf{\textcolor{CExample}{解题步骤：}}
	\begin{description}[style=nextline, leftmargin=2.8em, labelsep=0.5em, itemsep=0.45em]
		\item[\textbf{第一步（换元并定范围）：}] 令 $t=x+1$，由 $x>-1$ 得 $t>0$。通过除法变形：$y=2x+1+\frac{2}{x+1}$，代入 $t$ 得
			\[
				y = 2t+\frac{2}{t} -1,
				\qquad
				t>0.
			\]
			\par\textbf{理由：} 先换元并明确新元的正负情况，这是后续应用结论的前提.

		\item[\textbf{第二步（求对勾最小值）：}] 对于 $t>0$，$2t+\frac{2}{t}$ 的最小值由对勾结论给出：
			\[
				\begin{aligned}
					2t+\frac{2}{t} &\ge 2\sqrt{2t\cdot\frac{2}{t}} \\
					&= 2\sqrt{4} \\
					&= 4,
			\end{aligned}
			\]
			等号成立当且仅当 $2t=\frac{2}{t}$，即 $t=1$（$t>0$）.
			\par\textbf{理由：} 直接应用对勾函数 $x>0$ 时的最小值公式.

		\item[\textbf{第三步（回代得最小值）：}] 则
			\[
				y_{\min} = 3,
			\]
			此时 $t=1$，即 $x+1=1$，故 $x=0$.
			\par\textbf{理由：} 由 $y= (2t+\frac{2}{t}) -1$ 的最小值是对勾最小值减 $1$，即 $4-1=3$，回代自变量即得.
	\end{description}

	\textbf{\textcolor{CExample}{关键结论：}}
	\textbf{最小值为 $3$，此时 $x=0$}.
\end{examplebox}
